% TOP_PROBLEM: {"id": 30007156, "title": "Infinitely many prime geodesic images on spheres", "category_id": 6, "category": {"id": 6, "name": "geometry", "display_name": "Geometry", "description": "Euclidean and non-Euclidean geometry, geometric structures.", "slug": "geometry", "order_index": 6, "created_at": "2026-07-31T15:26:25.671Z"}, "status": "open"}
% FIRST_POSTED: 2026-10-06
% ATTEMPT_STATUS: solved
% OPENAI_CLAIM: {"schema_version":1,"source_repo":"https://github.com/openai/math","source_commit":"adc7f1241b42e322a6451854ab7e4b4c146bf78a","release_date":"2026-10-06","match":"full","resolution":"proved","independently_reviewed":false,"second_pass":"agreed","families":[{"family":"345","title":"Infinitely many closed geodesics on Riemannian spheres and closed three-manifolds","subject":"Differential geometry","manuscripts":[{"title":"Infinitely many closed geodesic images on every Riemannian sphere","pdf_path":"preprints/Infinitely-many-closed-geodesic-images-on-every-Riemannian-sphere-September-24-2026/paper.pdf","readme_path":"preprints/Infinitely-many-closed-geodesic-images-on-every-Riemannian-sphere-September-24-2026/README.md","date":"2026-09-24","theorem_ref":"thm:spheres (sections/introduction.tex:11); cor:three-manifolds (sections/introduction.tex:84); thm:doubling (sections/activation.tex:107)","statement_sources":[{"path":"preprints/Infinitely-many-closed-geodesic-images-on-every-Riemannian-sphere-September-24-2026/build/sections/introduction.tex","line":11,"label":"thm:spheres"},{"path":"preprints/Infinitely-many-closed-geodesic-images-on-every-Riemannian-sphere-September-24-2026/build/sections/introduction.tex","line":84,"label":"cor:three-manifolds"},{"path":"preprints/Infinitely-many-closed-geodesic-images-on-every-Riemannian-sphere-September-24-2026/build/sections/activation.tex","line":107,"label":"thm:doubling"}]}],"lean":null,"reasoning_trace":null}]}
\documentclass[11pt]{article}
\usepackage[margin=25mm]{geometry}
\usepackage{fontspec}
\setmainfont{Latin Modern Roman}
\usepackage{amsmath,amssymb,mathtools}
\usepackage{unicode-math}
\setmathfont{Latin Modern Math}
\usepackage{enumitem}
\usepackage{xurl,hyperref,bookmark}
\hypersetup{colorlinks=true,urlcolor=blue}
\setcounter{secnumdepth}{0}
\setlength{\parindent}{0pt}
\setlength{\parskip}{5pt}
\setlength{\emergencystretch}{3em}
\begin{document}
\section*{Infinitely many prime geodesic images on spheres}

\subsection{OpenAI's claim}
\textit{Infinitely many prime closed geodesics}: Every smooth Riemannian sphere of dimension at least two is claimed to have infinitely many prime closed geodesics with distinct images. The paper also obtains infinitude on every nonempty closed three-manifold.

\subsection{Mathematical statement of the claimed result}
\paragraph{Mathematical claim 1.} Whatever smooth metric is chosen on the standard sphere $S^n$, at each integer dimension $n\ge2$ the manuscript claims infinitely many prime closed geodesics having pairwise distinct images. Prime means not an iterate of a shorter closed geodesic. Its three-manifold corollary asserts the same conclusion for every nonempty closed smooth three-manifold and every smooth metric, allowing nonorientability.

\subsection{Outline of the argument}
\textit{Infinitely many prime closed geodesics}: Assuming finitely many prime images gives a uniform budget for local loop-space crossings, including degenerate root types. Reflection-equivariant cohomology and half-period translation compare ordinary and root-degree filtrations. Extra classes would force a quadratic discrepancy, but the exact-sequence and filling estimates allow only a linear budget; covering and three-manifold results provide the final transfer.

\subsection{Scope relative to this problem}
\textit{Infinitely many prime closed geodesics}: Every smooth Riemannian sphere of dimension at least two is claimed to have infinitely many prime closed geodesics with distinct images. The paper also obtains infinitude on every nonempty closed three-manifold. The parameters, regularity assumptions and normalization in this question are retained in applying that assertion.

With the hypotheses and conventions above, the claimed result addresses the whole mathematical assertion on this page.

\subsection{Formal verification}
\textit{Infinitely many prime closed geodesics}: The pinned catalogue lists no Lean formalization for this result.

This site has not run the Lean development. These links describe OpenAI’s selected formal statements and should be read with their stated scope.

\subsection{Status caveat}
OpenAI reports that an internal model produced the manuscripts. Their arguments have not been independently verified as proofs by this site. Any portion outside the documented formal scope may contain mistakes, and the catalogue label reports the announced mathematical claim.

\subsection{Sources}
\begin{itemize}
\item OpenAI, \emph{Infinitely many closed geodesic images on every Riemannian sphere}, 2026-09-24; citation key \texttt{OAI:Infinitely-many-closed-geodesic-images-on-every-Riemannian-sphere-September-24-2026}. \href{https://github.com/openai/math/blob/adc7f1241b42e322a6451854ab7e4b4c146bf78a/preprints/Infinitely-many-closed-geodesic-images-on-every-Riemannian-sphere-September-24-2026/paper.pdf}{PDF (GitHub)}; \href{https://github.com/openai/math/raw/adc7f1241b42e322a6451854ab7e4b4c146bf78a/preprints/Infinitely-many-closed-geodesic-images-on-every-Riemannian-sphere-September-24-2026/paper.pdf}{PDF (direct)}; \href{https://github.com/openai/math/blob/adc7f1241b42e322a6451854ab7e4b4c146bf78a/preprints/Infinitely-many-closed-geodesic-images-on-every-Riemannian-sphere-September-24-2026/README.md}{Manuscript page and citation}; \href{https://github.com/openai/math/blob/main/preprints/Infinitely-many-closed-geodesic-images-on-every-Riemannian-sphere-September-24-2026/paper.pdf}{Latest version}; Statement source: \href{https://github.com/openai/math/blob/adc7f1241b42e322a6451854ab7e4b4c146bf78a/preprints/Infinitely-many-closed-geodesic-images-on-every-Riemannian-sphere-September-24-2026/build/sections/introduction.tex\#L11}{thm:spheres}; Statement source: \href{https://github.com/openai/math/blob/adc7f1241b42e322a6451854ab7e4b4c146bf78a/preprints/Infinitely-many-closed-geodesic-images-on-every-Riemannian-sphere-September-24-2026/build/sections/introduction.tex\#L84}{cor:three-manifolds}; Statement source: \href{https://github.com/openai/math/blob/adc7f1241b42e322a6451854ab7e4b4c146bf78a/preprints/Infinitely-many-closed-geodesic-images-on-every-Riemannian-sphere-September-24-2026/build/sections/activation.tex\#L107}{thm:doubling}.
\item \href{https://github.com/openai/math/blob/adc7f1241b42e322a6451854ab7e4b4c146bf78a/CONTENTS.md}{OpenAI catalogue entry}.
\end{itemize}
\end{document}
